\documentclass[12pt]{article}

\usepackage{bbm}
\usepackage{fancyhdr}
\usepackage{titling}
\usepackage{enumitem}
\usepackage{amsmath}
\usepackage{dsfont}
\usepackage{amsfonts}
\usepackage{amssymb}
\usepackage{amsthm}
\usepackage{geometry}
\usepackage{tikz}
\usepackage{algorithm}
\usepackage{algpseudocode}
\geometry{verbose,tmargin=1in,bmargin=1in,lmargin=1in,rmargin=1in}
\renewcommand{\vec}[1]{\mathbf{#1}}
\newcommand{\dec}{\text{dec}}
\newcommand{\divi}{\mathbin{\textnormal{div}}}
\title{Homework 1}
\author{Satvik Prasad}
\date{}
\newtheorem{lemma}{Lemma}[section]
\fancypagestyle{plain}{%
  \fancyhf{}
  \fancyhead[L]{\theauthor}
  \fancyhead[R]{\thedate}
}
\pagestyle{plain}
\makeatletter
\def\@maketitle{%
  \newpage
  \null
  \vskip 1em%
  \begin{center}%
    \let \footnote \thanks
    {\LARGE \@title \par}%
    \vskip 1em%
  \end{center}%
  \par
\vskip 1em}
\makeatother
\usepackage{cmbright}

\newcommand{\mat}[1]{\mathbf{#1}}

% ---------- Probability theory ----------
% Core operators (optional argument-free; use as \E{X}, \Var{X}, \Cov{X}{Y})
\DeclareMathOperator{\EOp}{\mathbb{E}}
\DeclareMathOperator{\VarOp}{Var}
\DeclareMathOperator{\CovOp}{Cov}
\DeclareMathOperator{\CorrOp}{Corr}

\renewcommand{\P}[1]{\mathbb{P}\left(#1\right)}        % \P{A}
\newcommand{\E}[1]{\EOp\left[#1\right]}                % \E{X}
\newcommand{\Var}[1]{\VarOp\left(#1\right)}            % \Var{X}
\newcommand{\Cov}[2]{\CovOp\left(#1,\,#2\right)}       % \Cov{X}{Y}
\newcommand{\Corr}[2]{\CorrOp\left(#1,\,#2\right)}     % \Corr{X}{Y}
\newcommand{\SD}[1]{\operatorname{SD}\left(#1\right)}  % \SD{X}

% Subscripted versions, e.g. \Esub{\theta}{X}, \Psub{\theta}{A}
\newcommand{\Esub}[2]{\EOp_{#1}\left[#2\right]}
\newcommand{\Psub}[2]{\mathbb{P}_{#1}\left(#2\right)}

% Conditioning, indicators, independence
\newcommand{\given}{\,\middle|\,}                      % \E{X \given Y}
\newcommand{\ind}[1]{\mathbbm{1}\left\{#1\right\}}     % \ind{X > 0}
\newcommand{\indep}{\perp\!\!\!\perp}                  % X \indep Y
\newcommand{\iid}{\overset{\text{iid}}{\sim}}          % X_i \iid \Normal(0,1)

% Convergence
\newcommand{\convp}{\xrightarrow{\;p\;}}               % in probability
\newcommand{\convd}{\xrightarrow{\;d\;}}               % in distribution
\newcommand{\convas}{\xrightarrow{\text{a.s.}}}        % almost surely
\newcommand{\convLp}[1]{\xrightarrow{L^{#1}}}          % in L^p

% Probability spaces
\newcommand{\F}{\mathcal{F}}                           % sigma-algebra
\newcommand{\B}{\mathcal{B}}                           % Borel sigma-algebra
\newcommand{\pspace}{(\Omega, \F, \mathbb{P})}

% Common distributions
\newcommand{\Normal}{\mathcal{N}}
\newcommand{\Bern}{\operatorname{Bernoulli}}
\newcommand{\Bin}{\operatorname{Binomial}}
\newcommand{\Pois}{\operatorname{Poisson}}
\newcommand{\Unif}{\operatorname{Uniform}}
\newcommand{\Exp}{\operatorname{Exp}}
\newcommand{\Geom}{\operatorname{Geometric}}
\newcommand{\Gam}{\operatorname{Gamma}}
\newcommand{\Beta}{\operatorname{Beta}}

% Number sets
\newcommand{\R}{\mathbb{R}}
\newcommand{\N}{\mathbb{N}}
\newcommand{\Z}{\mathbb{Z}}


\begin{document}
\begin{enumerate}
  \item $\{q, s\} < \{u\} < \{r\} < \{f, h\} < \{k\} < \{g, m, t\} <
    \{w\} < \{p\}$

  \item
    \begin{enumerate}
      \item We prove that $\lg n! = \mathcal O(n \lg n)$. Notice that
        $$
        n! = 1 \cdot 2 \cdots n \leq n \cdot n \cdots n = n^n.
        $$
        Since $\lg x$ is monotonically increasing, we can conclude that
        $$
        \lg n! \leq \lg (n^{n}) = n \lg n.
        $$
        Hence, for $n \geq n_0 = 1$, there exists a $c = 1$ such that
        $$
        \lg n! \leq c \cdot n \lg n,
        $$
        meaning $\lg n! = \mathcal O(n \lg n)$.

      \item Notice that for all $n \geq 2$,
        $$
        n! = \big(1 \cdot 2 \cdots \lfloor n/2 \rfloor\big)
        \big((\lfloor n/2 \rfloor + 1) \cdots n\big).
        $$
        Hence,
        $$
        n! \geq \underbrace{\lfloor n/2\rfloor!}_{\geq 1}
        \times \lfloor n/2 \rfloor^{\,n - \lfloor n/2 \rfloor}.
        $$
        Since $n - \lfloor n/2 \rfloor \geq n/2 \geq 1$, we can conclude
        $$
        n! \geq \lfloor n/2 \rfloor^{n/2}.
        $$
        Finally, since for all $n \geq 2$, $\lfloor n/2 \rfloor \geq
        n/4$, we get
        \begin{align*}
          \lg n! &\geq \frac{n}{2} \lg \lfloor n/2 \rfloor \\
          &\geq \frac{n}{2} \lg \frac{n}{4} \\
          &= \frac{n}{2} \lg n - n.
        \end{align*}
        Thus, for $n \geq 16 = 2^4$, we get
        $$
        \frac{1}{4} \lg n - 1 \geq 0 \;\Rightarrow\; \frac{n}{4} \lg
        n - n \geq 0,
        $$
        meaning $\frac{n}{2} \lg n - n \geq \frac{n}{4} \lg n$. Thus, for
        $n \geq n_0 = 16$, there exists a $c = 1/4$ such that
        $$
        \lg n! \geq c \, n \lg n,
        $$
        meaning $\lg n! = \Omega(n \lg n)$.

      \item Note that for $a, b > 1$, $\log_a n = \lg n / \lg a$. Thus, for all
        $n \geq 1$, there exists a $c = \dfrac{\lg b}{\lg a}$ such that
        $$
        \log_a n = \frac{\lg n}{\lg a}
        = \frac{\lg b}{\lg a} \cdot \frac{\lg n}{\lg b}
        = c \cdot \log_b n
        \leq c \cdot \log_b n.
        $$
        Thus, $\log_a n = \mathcal O(\log_b n)$.
    \end{enumerate}

  \item
    \begin{enumerate}
      \item \textsc{BinarySearch} satisfies $T(1) = \mathcal O(1)$ and
        $T(n) = T(n/2) + \mathcal O(1)$, meaning $T(n) = \mathcal O(\lg n)$.

      \item \textsc{Maximum} satisfies $T(1) = \mathcal O(1)$ and
        $T(n) = 2T(n/2) + \mathcal O(1)$, meaning $T(n) = \mathcal O(n)$.

      \item \textsc{AlternateMaximum} satisfies $T(1) = \mathcal O(1)$ and
        $T(n) = T(n-1) + \mathcal O(1)$, meaning $T(n) = \mathcal O(n)$.

      \item \textsc{MyAlgo} satisfies $T(n) = \mathcal O(1)$ for all
        $n \in \{1, \ldots, 10\}$, and
        $$
        T(n) = 2T\!\left(\tfrac{3}{4}n\right) + \mathcal O(n).
        $$
        Using the master theorem, since $\log_{4/3} 2 > 1$, we get
        $$
        T(n) = \mathcal O\!\left(n^{\log_{4/3} 2}\right).
        $$
    \end{enumerate}

  \item Since $n$ can be halved at most $\lg(n/k)$ times before it is less
    than $k$, the recursion can be viewed as a binary tree of height
    $\lg(n/k)$, where the work at a leaf is $\mathcal O(k \lg k)$, and the work
    at an internal node is $\mathcal O(k)$ more than the work at its children.

    Let $L = \lg(n/k)$, and for fixed $n$ let $R(h)$ be the total work
    processing up to height $h$ in the tree, so that $T(n) = R(L)$. From the
    base case,
    $$
    R(0) = \mathcal O\big(k \lg k \cdot (n/k)\big) = \mathcal O(n \lg k),
    $$
    since there are $n/k$ leaves, each requiring $\mathcal O(k \lg k)$ work.
    Additionally, since height $h$ has $(n/k) \cdot 2^{-h}$ nodes, each
    requiring $\mathcal O(k)$ additional work,
    \begin{align*}
      R(h) &= R(h-1) + \mathcal O\big((n/k) \cdot 2^{-h} \cdot k\big)
      = R(h-1) + \mathcal O(2^{-h} n), \qquad h \geq 1, \\
      R(h) &= \mathcal O(n \lg k) + \sum_{j=1}^{h} \mathcal O(2^{-j} n).
    \end{align*}
    Hence, using properties of the geometric series,
    $$
    T(n) = R(L)
    = \mathcal O(n \lg k) + \frac{n}{2} \sum_{j=0}^{L-1} \mathcal O(2^{-j})
    = \mathcal O(n \lg k) + \mathcal O(n - k)
    = \mathcal O(n \lg k),
    $$
    as required.
\end{enumerate}
\end{document}
